\chapter{2009年大纲II卷（文）}

\section{单选题}

\begin{problem}
已知全集 \(U = \{1, 2, 3, 4, 5, 6, 7, 8\}\)，\(M = \{1, 3, 5, 7\}\)，\(N = \{5, 6, 7\}\)，则 \(\complement_U(M \cup N) =\)（\quad）

\choices
    {\( \{5, 7\} \)}
    {\(\{2, 4\}\)}
    {\(\{2, 4, 8\}\)}
    {\(\{1,3,5,6,7\}\)}
\end{problem}

\begin{answer}
C
\end{answer}

\begin{solution}
由 \(M\cup N=\{1,3,5,6,7\}\)，得
\(\complement_U(M\cup N)=\{2,4,8\}\)，所以选 C.
\end{solution}

\begin{problem}
函数 \( y = \sqrt{-x} \)（\(x \leqslant 0\)）的反函数是（\quad）

\choices
    {\(y = x^{2}\)（\(x\geqslant 0\)）}
    {\(y = -x^{2}\)（\(x\geqslant 0\)）}
    {\(y = x^{2}\)（\(x\leqslant 0\)）}
    {\(y = -x^{2}\)（\(x\leqslant 0\)）}
\end{problem}

\begin{answer}
B
\end{answer}

\begin{solution}
由 \(y=\sqrt{-x}\) 可知 \(y\geqslant0\). 交换 \(x\)，\(y\) 得 \(x=\sqrt{-y}\)，因而 \(x\geqslant0\) 且 \(y=-x^2\). 所以反函数为 \(y=-x^2\ (x\geqslant0)\)，选 B.
\end{solution}

\begin{problem}
函数 \( y = \log_2\frac{2 - x}{2 + x} \) 的图象（\quad）

\choices
    {关于原点对称}
    {关于直线 \( y = -x \) 对称}
    {关于 \(y\) 轴对称}
    {关于直线 \( y = x \) 对称}
\end{problem}

\begin{answer}
A
\end{answer}

\begin{solution}
定义域为 \((-2,2)\). 对任意 \(x\in(-2,2)\)，有 \(f(-x)=\log_2\dfrac{2+x}{2-x}=-\log_2\dfrac{2-x}{2+x}=-f(x)\)，故函数为奇函数，图象关于原点对称，选 A.
\end{solution}

\begin{problem}
已知 \(\triangle ABC\) 中，\(\cot A = -\frac{12}{5}\)，则 \(\cos A =\)（\quad）

\choices
    {\(\frac{12}{13}\)}
    {\(\frac{5}{13}\)}
    {\(-\frac{5}{13}\)}
    {\(-\frac{12}{13}\)}
\end{problem}

\begin{answer}
D
\end{answer}

\begin{solution}
三角形内角满足 \(0<A<\pi\)，且 \(\cot A<0\)，所以 \(A\) 为钝角，\(\cos A<0\). 由 \(1+\cot^2A=\csc^2A\) 或 \(\sin^2A+\cos^2A=1\)，得
\(\cos^2A=\dfrac{\cot^2A}{1+\cot^2A}=\dfrac{144}{169}\). 因此 \(\cos A=-\dfrac{12}{13}\)，选 D.
\end{solution}

\begin{problem}
已知正四棱柱 \(ABCD - A_{1}B_{1}C_{1}D_{1}\) 中，\(AA_{1} = 2AB\)，\(E\) 为 \(AA_{1}\) 中点，则异面直线 \(BE\) 与 \(CD_{1}\) 所成的角的余弦值为（\quad）

\choices
    {\(\frac{\sqrt{10}}{10}\)}
    {\(\frac{1}{5}\)}
    {\(\frac{3\sqrt{10}}{10}\)}
    {\(\frac{3}{5}\)}
\end{problem}

\begin{answer}
C
\end{answer}

\begin{solution}
令 \(AB=1\)，建立直角坐标系：\(A(0,0,0)\)，\(B(1,0,0)\)，\(C(1,1,0)\)，\(D(0,1,0)\)，\(A_1(0,0,2)\)，则 \(E(0,0,1)\)，\(D_1(0,1,2)\). 取方向向量 \(\overrightarrow{BE}=(-1,0,1)\)，\(\overrightarrow{CD_1}=(-1,0,2)\)，则
\(\overrightarrow{BE}\cdot\overrightarrow{CD_1}=3\)，\(|\overrightarrow{BE}|=\sqrt2\)，\(|\overrightarrow{CD_1}|=\sqrt5\). 因而两异面直线所成的角的余弦为 \(\dfrac{3}{\sqrt{10}}=\dfrac{3\sqrt{10}}{10}\)，选 C.
\end{solution}

\begin{problem}
已知向量 \(\bs{a} = (2,1)\)，\(\bs{a} \cdot \bs{b} = 10\)，\(|\bs{a} + \bs{b}| = 5\sqrt{2}\)，则 \(|\bs{b}| =\)（\quad）

\choices
    {\(\sqrt{5}\)}
    {\(\sqrt{10}\)}
    {\(5\)}
    {\(25\)}
\end{problem}

\begin{answer}
C
\end{answer}

\begin{solution}
由 \(|\bs{a}|^2=2^2+1^2=5\) 及平方和公式，得
\(|\bs{a}+\bs{b}|^2=|\bs{a}|^2+2\bs{a}\cdot\bs{b}+|\bs{b}|^2\). 代入已知量，\(50=5+2\times10+|\bs{b}|^2\)，所以 \(|\bs{b}|^2=25\)，从而 \(|\bs{b}|=5\)，选 C.
\end{solution}

\begin{problem}
设 \(a = \lg \e\)，\(b = (\lg \e)^2\)，\(c = \lg \sqrt{\e}\)，则（\quad）

\choices
    {\(a > b > c\)}
    {\(a > c > b\)}
    {\(c > a > b\)}
    {\(c > b > a\)}
\end{problem}

\begin{answer}
B
\end{answer}

\begin{solution}
令 \(t=\lg\e=\dfrac1{\ln10}\). 因为 \(\ln10>2\)，所以 \(0<t<\dfrac12\). 又 \(b=t^2\)，\(c=\lg\sqrt{\e}=\dfrac12\lg\e=\dfrac t2\)，故 \(t>\dfrac t2>t^2\)，即 \(a>c>b\)，选 B.
\end{solution}

\begin{problem}
双曲线 \(\frac{x^2}{6} - \frac{y^2}{3} = 1\) 的渐近线与圆 \((x - 3)^2 + y^2 = r^2\)（\(r > 0\)）相切，则 \(r =\)（\quad）

\choices
    {\(\sqrt{3}\)}
    {\(2\)}
    {\(3\)}
    {\(6\)}
\end{problem}

\begin{answer}
A
\end{answer}

\begin{solution}
双曲线的渐近线为 \(y=\pm\dfrac1{\sqrt2}x\)，即 \(x\mp\sqrt2y=0\). 圆心为 \(P(3,0)\)，圆与渐近线相切，所以半径等于圆心到渐近线的距离：
\(r=\dfrac{|3|}{\sqrt{1+2}}=\sqrt3\). 因此选 A.
\end{solution}

\begin{problem}
若将函数 \( y = \tan \left( \omega x + \frac{\pi}{4} \right) \) （\(\omega > 0\)）的图象向右平移 \( \frac{\pi}{6} \) 个单位长度后，与函数 \( y = \tan \left( \omega x + \frac{\pi}{6} \right) \) 的图象重合，则 \( \omega \) 的最小值为（\quad）

\choices
    {\(\frac{1}{6}\)}
    {\(\frac{1}{4}\)}
    {\(\frac{1}{3}\)}
    {\(\frac{1}{2}\)}
\end{problem}

\begin{answer}
D
\end{answer}

\begin{solution}
向右平移 \(\dfrac\pi6\) 后，原图象变为
\(y=\tan\left(\omega x-\dfrac{\omega\pi}{6}+\dfrac\pi4\right)\). 两正切函数图象重合，当且仅当两相位相差 \(k\pi\)，故
\(-\dfrac{\omega\pi}{6}+\dfrac\pi4=\dfrac\pi6+k\pi\)，从而 \(\omega=\dfrac12-6k\). 由 \(\omega>0\) 及 \(k\in\mathbb Z\)，最小值在 \(k=0\) 时取得，为 \(\dfrac12\)，选 D.
\end{solution}

\begin{problem}
甲，乙两人从 \(4\text{ 门}\)课程中各选修 \(2\text{ 门}\)，则甲，乙所选的课程中恰有 \(1\text{ 门}\)相同的选法有（\quad）

\choices
    {\(6\text{ 种}\)}
    {\(12\text{ 种}\)}
    {\(24\text{ 种}\)}
    {\(30\text{ 种}\)}
\end{problem}

\begin{answer}
C
\end{answer}

\begin{solution}
先从 \(4\) 门课程中为甲选取 \(2\) 门，有 \(\mathrm C_4^2\) 种. 固定甲的选择后，乙要恰好选中其中 \(1\) 门，需从甲所选的 \(2\) 门中选 \(1\) 门，再从其余 \(2\) 门中选 \(1\) 门，有 \(\mathrm C_2^1\mathrm C_2^1\) 种. 因而总数为 \(\mathrm C_4^2\mathrm C_2^1\mathrm C_2^1=6\times2\times2=24\)，选 C.
\end{solution}

\begin{problem}
已知直线 \( y = k(x + 2) \) （\( k > 0 \)）与抛物线 \( C: y^2 = 8x \) 相交于 \(A\)，\(B\) 两点，\(F\) 为 \(C\) 的焦点. 若 \( |FA| = 2|FB| \)，则 \( k = \)（\quad）

\choices
    {\( \frac{1}{3} \)}
    {\( \frac{\sqrt{2}}{3} \)}
    {\( \frac{2}{3} \)}
    {\( \frac{2\sqrt{2}}{3} \)}
\end{problem}

\begin{answer}
D
\end{answer}

\begin{solution}
抛物线 \(y^2=8x\) 的参数方程可写为 \(x=2t^2\)，\(y=4t\)，其焦点为 \(F(2,0)\). 交点对应参数 \(t_1\)，\(t_2\)，代入直线 \(y=k(x+2)\) 得 \(2t=k(t^2+1)\)，即 \(kt^2-2t+k=0\). 因而 \(t_1t_2=1\). 对抛物线上的参数点，有
\(|F(2t^2,4t)|=\sqrt{(2t^2-2)^2+(4t)^2}=2(t^2+1)\). 所以
\(\dfrac{|FA|}{|FB|}=\dfrac{t_1^2+1}{t_2^2+1}=t_1^2\)（交换 \(A\)，\(B\) 时取倒数），由已知比例得 \(\{t_1^2,t_2^2\}=\{2,\dfrac12\}\). 因为 \(t_1t_2=1>0\)，两个参数同号；又由韦达定理 \(t_1+t_2=\dfrac2k>0\)，故 \(t_1\)，\(t_2>0\)，从而 \(t_1+t_2=\sqrt2+\dfrac1{\sqrt2}=\dfrac3{\sqrt2}\). 所以 \(k=\dfrac{2\sqrt2}{3}\)，选 D.
\end{solution}

\begin{problem}
纸制的正方体的六个面根据其方位分别标记为上，下，东，南，西，北. 现有沿该正方体的一些棱将正方体剪开，外面朝上展平，得到如图的平面图形，则标“\(\triangle\)”的面的方位是（\quad）

\FigureLayoutDeclare{img/2009/syllabus_paper_2_liberal/q12_fig1.png}{6.0cm}{4bp 3bp 4bp 4bp}
\bitmapfigure[max width=0.42\linewidth]{img/2009/syllabus_paper_2_liberal/q12_fig1.png}

\choices
    {南}
    {北}
    {西}
    {下}
\end{problem}

\begin{answer}
B
\end{answer}

\begin{solution}
将图中水平一排的四个面折起后，它们依次组成正方体的侧面. 由图中相邻标有“上”和“东”的两个面，可确定从左到右四个面依次为“下”“西”“上”“东”. 标有三角形的面与最左面的“下”面相邻，沿该棱折起后位于“北”方位. 因此选 B.
\end{solution}

\section{填空题}

\begin{problem}
设等比数列 \(\{a_{n}\}\) 的前 \(n\) 项和为 \(S_{n}\)，若 \(a_{1} = 1\)，\(S_{6} = 4S_{3}\)，则 \(a_{4} = \)\fillinblank{}.
\end{problem}

\begin{answer}
\(3\)
\end{answer}

\begin{solution}
设等比数列的公比为 \(q\). 当 \(q=1\) 时，\(S_6=6\)，而 \(4S_3=12\)，不满足条件，故 \(q\neq1\). 又 \(S_6=(1+q+q^2)(1+q^3)\)，且 \(S_3=1+q+q^2>0\)，所以 \(1+q^3=4\)，即 \(q^3=3\). 因而 \(a_4=a_1q^3=3\).
\end{solution}

\begin{problem}
\(\left(x\sqrt{y} - y\sqrt{x}\right)^4\) 的展开式中 \(x^3 y^3\) 的系数为\fillinblank{}.
\end{problem}

\begin{answer}
\(6\)
\end{answer}

\begin{solution}
展开通项为 \(\mathrm C_4^k\left(x\sqrt y\right)^{4-k}\left(-y\sqrt x\right)^k\). 该项中 \(x\)，\(y\) 的幂次分别为 \(4-\dfrac k2\)，\(2+\dfrac k2\). 要得到 \(x^3y^3\)，必须有 \(4-\dfrac k2=3\)，即 \(k=2\). 因此所求系数为 \(\mathrm C_4^2(-1)^2=6\).
\end{solution}

\begin{problem}
已知圆 \(O: x^{2} + y^{2} = 5\) 和点 \(A(1,2)\)，则过 \(A\) 且与圆 \(O\) 相切的直线与两坐标轴围成的三角形的面积等于\fillinblank{}.
\end{problem}

\begin{answer}
\(\frac{25}{4}\)
\end{answer}

\begin{solution}
点 \(A(1,2)\) 满足 \(1^2+2^2=5\)，所以 \(A\) 在圆上. 半径 \(OA\) 的斜率为 \(2\)，圆在 \(A\) 处的切线斜率为 \(-\dfrac12\)，故切线方程为 \(y-2=-\dfrac12(x-1)\)，即 \(x+2y-5=0\). 它与 \(x\) 轴、\(y\) 轴的截距分别为 \(5\)，\(\dfrac52\)，所以所围直角三角形的面积为 \(\dfrac12\times5\times\dfrac52=\dfrac{25}{4}\).
\end{solution}

\begin{problem}
设 \(OA\) 是球 \(O\) 的半径，\(M\) 是 \(OA\) 的中点，过 \(M\) 且与 \(OA\) 成 \(45^{\circ}\) 角的平面截球 \(O\) 的表面得到圆 \(C\). 若圆 \(C\) 的面积等于 \(\frac{7\pi}{4}\)，则球 \(O\) 的表面积等于\fillinblank{}.
\end{problem}

\begin{answer}
\(8\pi\)
\end{answer}

\begin{solution}
设球半径为 \(R\)，则 \(OM=\dfrac R2\). 平面与 \(OA\) 的夹角为 \(45^{\circ}\)，所以球心到截面的距离为 \(d=OM\sin45^{\circ}=\dfrac{R}{2\sqrt2}\). 设截面圆半径为 \(r\)，则 \(r^2=R^2-d^2=R^2-\dfrac{R^2}{8}=\dfrac{7R^2}{8}\). 由截面圆面积 \(\pi r^2=\dfrac{7\pi}{4}\)，得 \(R^2=2\). 因此球的表面积为 \(4\pi R^2=8\pi\).
\end{solution}

\section{解答题}

\begin{problem}
已知等差数列 \(\{a_{n}\}\) 中，\(a_{3}a_{7} = -16\)，\(a_{4} + a_{6} = 0\). 求 \(\{a_{n}\}\) 的前 \(n\) 项和 \(S_{n}\).
\end{problem}

\begin{solution}
设等差数列的首项为 \(a_1\)，公差为 \(d\). 由
\(a_4+a_6=0\)，得
\[
(a_1+3d)+(a_1+5d)=0
\]
即 \(a_1=-4d\). 因而
\(a_3=a_1+2d=-2d\)，\(a_7=a_1+6d=2d\).
代入 \(a_3a_7=-16\)，得
\(-4d^2=-16\)，所以 \(d=\pm2\).
又
\[
S_n=\frac n2\bigl(2a_1+(n-1)d\bigr)
    =\frac{nd(n-9)}2
\]
当 \(d=2\) 时，\(S_n=n(n-9)\)；当 \(d=-2\) 时，
\(S_n=-n(n-9)\).
\end{solution}

\begin{problem}
设 \(\triangle ABC\) 的内角 \(A\)，\(B\)，\(C\) 的对边长分别为 \(a\)，\(b\)，\(c\)，\(\cos (A - C) + \cos B = \frac{3}{2}\)，\(b^2 = ac\)，求 \(B\).
\end{problem}

\begin{solution}
由 \(A+B+C=\pi\)，有
\[
\cos B=\cos\bigl(\pi-(A+C)\bigr)=-\cos(A+C)
\]
因此
\[
\cos(A-C)+\cos B
=\cos(A-C)-\cos(A+C)
=2\sin A\sin C=\frac32
\]
即 \(\sin A\sin C=\dfrac34\).
由正弦定理，\(a=2R\sin A\)，\(b=2R\sin B\)，
\(c=2R\sin C\). 条件 \(b^2=ac\) 给出
\[
\sin^2B=\sin A\sin C=\frac34
\]
所以 \(B=\dfrac{\pi}{3}\) 或 \(B=\dfrac{2\pi}{3}\).
若 \(B=\dfrac{2\pi}{3}\)，则 \(A+C=\dfrac{\pi}{3}\)，而
\[
\sin A\sin C\leqslant
\sin^2\frac{A+C}{2}=\sin^2\frac{\pi}{6}=\frac14
\]
与 \(\sin A\sin C=\dfrac34\) 矛盾. 故
\(B=\dfrac{\pi}{3}\).
\end{solution}

\begin{problem}
如图，直三棱柱 \(ABC - A_{1}B_{1}C_{1}\) 中，\(AB \perp AC\)，\(D\)，\(E\) 分别为 \(A A_{1}\)，\(B_{1}C\) 的中点，\(DE \perp\) 平面 \(BCC_{1}\).

\begin{enumerate}
\item 证明： \(AB = AC\)；
\item 设二面角 \(A - BD - C\) 为 \(60^{\circ}\)，求 \(B_{1}C\) 与平面 \(BCD\) 所成的角的大小.
\end{enumerate}

\bitmapfigure[max width=0.46\linewidth]{img/2009/syllabus_paper_2_liberal/q19_fig1.png}
\end{problem}

\begin{solution}
\begin{enumerate}
\item 设 \(AB=b\)，\(AC=c\)，\(AA_1=h\)，建立直角坐标系：
\(A(0,0,0)\)，\(B(b,0,0)\)，\(C(0,c,0)\)，\(A_1(0,0,h)\).
则
\(D\left(0,0,\frac h2\right)\)，\(E\left(\frac b2,\frac c2,\frac h2\right)\).
平面 \(BCC_1\) 的一个法向量为 \((c,b,0)\)，而
\(\overrightarrow{DE}=\left(\dfrac b2,\dfrac c2,0\right)\).
由 \(DE\perp\) 平面 \(BCC_1\)，两向量平行，故
\[
\frac b2:\frac c2=c:b
\]
于是 \(b^2=c^2\). 又 \(b\)，\(c>0\)，所以 \(b=c\)，即
\(AB=AC\).
\item 由第一问可令 \(AB=AC=1\). 此时取
\(A(0,0,0)\)，\(B(1,0,0)\)，\(C(0,1,0)\)，\(D\left(0,0,\frac h2\right)\)，\(B_1(1,0,h)\).
平面 \(ABD\) 的法向量可取 \(\bs{n}_1=(0,1,0)\)，
平面 \(CBD\) 的法向量可取
\(\bs{n}_2=\left(\dfrac h2,\dfrac h2,1\right)\).
由二面角 \(A-BD-C=60^\circ\)，得
\[
\frac{\bs{n}_1\cdot\bs{n}_2}
{|\bs{n}_1||\bs{n}_2|}
=\frac{h/2}{\sqrt{h^2/2+1}}=\cos60^\circ=\frac12
\]
解得 \(h^2=2\)，故 \(h=\sqrt2\).
取 \(\overrightarrow{B_1C}=(1,-1,\sqrt2)\)，
平面 \(BCD\) 的法向量可取
\(\bs{n}=(1,1,\sqrt2)\). 设直线 \(B_1C\) 与平面
\(BCD\) 所成的角为 \(\theta\)，则
\[
\sin\theta
=\frac{|\overrightarrow{B_1C}\cdot\bs{n}|}
{|\overrightarrow{B_1C}||\bs{n}|}
=\frac2{2\cdot2}=\frac12
\]
所以 \(\theta=30^\circ\).
\end{enumerate}
\end{solution}

\begin{problem}
某车间甲组有 \(10\text{ 名}\)工人，其中有 \(4\text{ 名}\)女工人；乙组有 \(10\text{ 名}\)工人，其中有 \(6\text{ 名}\)女工人. 现采用分层抽样方法（层内采用不放回简单随机抽样）从甲，乙两组中共抽取 \(4\text{ 名}\)工人进行技术考核.

\begin{enumerate}
\item 求从甲，乙两组各抽取的人数；
\item 求从甲组抽取的工人中恰有 \(1\text{ 名}\)女工人的概率；
\item 求抽取的 \(4\text{ 名}\)工人中恰有 \(2\text{ 名}\)男工人的概率.
\end{enumerate}
\end{problem}

\begin{solution}
\begin{enumerate}
\item 甲、乙两组人数相同，按比例分层抽样时各抽取
\(4\times\dfrac{10}{20}=2\) 名工人.
\item 甲组有 \(4\) 名女工、\(6\) 名男工. 从甲组抽取
\(2\) 人的总数为 \(\mathrm C_{10}^2\). 恰有 \(1\) 名女工的取法数为
\(\mathrm C_4^1\mathrm C_6^1\)，所以所求概率为
\[
\frac{\mathrm C_4^1\mathrm C_6^1}{\mathrm C_{10}^2}
=\frac{24}{45}=\frac8{15}
\]
\item 每组都抽取 \(2\) 人，总取法数为
\(\left(\mathrm C_{10}^2\right)^2\). 恰有 \(2\) 名男工时，甲组抽取的男工人数
\(i\) 只能为 \(0\)，\(1\)，\(2\)，相应取法数为
\[
\mathrm C_6^0\mathrm C_4^2\mathrm C_4^2\mathrm C_6^0
+\mathrm C_6^1\mathrm C_4^1\mathrm C_4^1\mathrm C_6^1
+\mathrm C_6^2\mathrm C_4^0\mathrm C_4^0\mathrm C_6^2
=36+576+225=837
\]
故所求概率为
\[
\frac{837}{\left(\mathrm C_{10}^2\right)^2}
=\frac{837}{2025}=\frac{31}{75}
\]
\end{enumerate}
\end{solution}

\begin{problem}
设函数 \( f(x) = \frac{1}{3} x^3 - (1 + a)x^2 + 4ax + 24a \)，其中常数 \( a > 1 \).

\begin{enumerate}
\item 讨论 \( f(x) \) 的单调性；
\item 若当 \(x \geqslant 0\) 时，\(f(x) > 0\) 恒成立，求 \(a\) 的取值范围.
\end{enumerate}
\end{problem}

\begin{solution}
\begin{enumerate}
\item 求导得
\[
f'(x)=x^2-2(1+a)x+4a=(x-2)(x-2a)
\]
由于 \(a>1\)，有 \(2<2a\)，故
\(f'(x)>0\) 在 \((-\infty,2)\) 和 \((2a,+\infty)\) 上成立，
\(f'(x)<0\) 在 \((2,2a)\) 上成立. 所以 \(f(x)\) 在
\((-\infty,2)\) 和 \((2a,+\infty)\) 上单调递增，在
\((2,2a)\) 上单调递减.
\item 当 \(x\geqslant0\) 时，\(f\) 在 \([0,2]\) 上递增，在
\([2,2a]\) 上递减，在 \([2a,+\infty)\) 上递增. 因此只需比较
\(f(0)\) 与 \(f(2a)\). 有
\[
f(0)=24a>0
\]
且
\[
f(2a)=\frac43a(6-a)(a+3)
\]
因为 \(a>1\)，所以 \(a>0\) 且 \(a+3>0\). 要使
\(f(x)>0\) 对所有 \(x\geqslant0\) 成立，必须且只需
\(f(2a)>0\)，即 \(a<6\). 故 \(a\) 的取值范围为
\(1<a<6\).
\end{enumerate}
\end{solution}

\begin{problem}
已知椭圆 \(C: \frac{x^2}{a^2} + \frac{y^2}{b^2} = 1 \) （\(a > b > 0\)）的离心率为 \(\frac{\sqrt{3}}{3}\)，过右焦点 \(F\) 的直线 \(l\) 与 \(C\) 相交于 \(A\)，\(B\) 两点，当 \(l\) 的斜率为 \(1\) 时，坐标原点 \(O\) 到 \(l\) 的距离为 \(\frac{\sqrt{2}}{2}\).

\begin{enumerate}
\item 求 \(a\)，\(b\) 的值；
\item \(C\) 上是否存在点 \(P\)，使得当 \(l\) 绕 \(F\) 转到某一位置时，有 \(\overrightarrow{OP} = \overrightarrow{OA} + \overrightarrow{OB}\) 成立？若存在，求出所有的 \(P\) 的坐标与 \(l\) 的方程；若不存在，说明理由.
\end{enumerate}
\end{problem}

\begin{solution}
\begin{enumerate}
\item 设椭圆焦距的一半为 \(c\). 由离心率
\(\dfrac ca=\dfrac{\sqrt3}{3}\)，得
\(c=\dfrac{\sqrt3}{3}a\)，从而
\[
b^2=a^2-c^2=\frac23a^2
\]
右焦点为 \(F(c,0)\). 当直线 \(l\) 的斜率为 \(1\) 时，
\(l\) 的方程为 \(y=x-c\)，即 \(x-y-c=0\). 由原点到
\(l\) 的距离为 \(\dfrac{\sqrt2}{2}\)，得
\[
\frac c{\sqrt2}=\frac{\sqrt2}{2}
\]
所以 \(c=1\). 于是 \(a=\sqrt3\)，\(b=\sqrt2\).
\item 设过 \(F(1,0)\) 的非竖直直线为
\[
l:y=m(x-1)
\]
将其代入椭圆方程
\(\dfrac{x^2}{3}+\dfrac{y^2}{2}=1\)，得
\[
(2+3m^2)x^2-6m^2x+3m^2-6=0
\]
设交点 \(A\)，\(B\) 的横坐标为 \(x_1\)，\(x_2\)，则
\(x_1+x_2=\frac{6m^2}{2+3m^2}\)，\(y_1+y_2=m(x_1+x_2-2) =-\frac{4m}{2+3m^2}\).
若 \(P=(X,Y)\) 满足
\(\overrightarrow{OP}=\overrightarrow{OA}+\overrightarrow{OB}\)，
则
\(X=\frac{6m^2}{2+3m^2}\)，\(Y=-\frac{4m}{2+3m^2}\).
再由 \(P\) 在椭圆上，
\[
\frac{X^2}{3}+\frac{Y^2}{2}
=\frac{4m^2}{2+3m^2}=1
\]
所以 \(m^2=2\). 当 \(m=\sqrt2\) 时，
\(P\left(\frac32,-\frac{\sqrt2}{2}\right)\)，\(l:y=\sqrt2(x-1)\)
即 \(l:x=\dfrac{\sqrt2}{2}y+1\). 当 \(m=-\sqrt2\) 时，
\(P\left(\frac32,\frac{\sqrt2}{2}\right)\)，\(l:y=-\sqrt2(x-1)\)
即 \(l:x=-\dfrac{\sqrt2}{2}y+1\).
最后，竖直直线 \(x=1\) 与椭圆交于两点时，其横坐标和为
\(2\)，纵坐标和为 \(0\)，对应 \(P=(2,0)\)，而
\((2,0)\) 不在椭圆上，故没有遗漏.
\end{enumerate}
\end{solution}
